\documentclass[aspectratio=169,9pt,demo]{beamer}
\usepackage[utf8]{inputenc}
\usepackage{mathtools}
\usepackage{multicol}
\usepackage{mathptmx}
\usepackage{multirow}
\usepackage{pseudocode}
\usepackage{caption}
\usepackage{subcaption} 
\usepackage[table]{xcolor}
\usepackage{colortbl}
\usepackage{tikz}
\usetikzlibrary{shapes.geometric, arrows}
\usepackage{pgfplots}
\pgfplotsset{compat=1.18}
\usepackage{pgfplotstable}
\usepackage{setspace} 

\usetheme{CambridgeUS}
\setbeamertemplate{headline}{%
  \leavevmode%
  \hbox{%
    \begin{beamercolorbox}[wd=0.55\paperwidth,ht=0.3cm,dp=0.2cm,center]{section in head/foot}
      \textbf{\insertsection} 
    \end{beamercolorbox}%
    \begin{beamercolorbox}[wd=0.45\paperwidth,ht=0.3cm,dp=0.2cm,center]{subsection in head/foot}
      \textbf{\insertsubsection} 
    \end{beamercolorbox}
  }%
}

\usecolortheme{orchid}

\definecolor{myNewColorA}{RGB}{39,52,139}
\definecolor{myNewcolorA}{RGB}{226,182,0}

\setbeamercolor{block title}{bg=myNewColorA,fg=white}
\setbeamercolor{block body}{bg=myNewColorA!20,fg=black}
\setbeamercolor{block title alerted}{bg=white, fg=myNewColorA}
\setbeamercolor{block body alerted}{bg=black!20, fg=white}

\setbeamercolor*{block title example}{bg=myNewColorA, fg = white}
\setbeamercolor*{block body example}{bg=myNewColorA!20, fg = black}
\usebeamercolor[myNewColorA]{block title alerted}
\setbeamercolor*{palette primary}{bg=myNewcolorA, fg = black}
\setbeamercolor*{palette secondary}{bg=myNewcolorA, fg = black}
\setbeamercolor*{palette tertiary}{bg=myNewcolorA, fg = white}
\setbeamercolor*{titlelike}{fg=myNewColorA}
\setbeamercolor*{title}{bg=myNewColorA, fg = white}
\setbeamercolor*{item}{fg=myNewColorA}
\setbeamercolor*{caption name}{fg=myNewcolorA}
\usefonttheme{professionalfonts}
\usepackage{natbib}
\usepackage{hyperref}
\setbeamersize{text margin left = 2mm, text margin right = 2mm}

%------------------------------------------------------------
\titlegraphic{\includegraphics[height=1.5cm]{ressources/images/logo.jpg}} 

\setbeamerfont{title}{size=\large}
\setbeamerfont{author}{size=\normalsize}
\setbeamerfont{date}{size=\small}
\setbeamerfont{institute}{size=\small}

\title[Mémoire de Licence Professionnelle]{PAYQR - APPLICATION DE PAIEMENT QR-CODE}

% Configuration pour Beamer
\author[NJANKO, ETOA, ZIEM, NKONGMENOCK \& MARTIAL]{%
  Présenté par : \\
  ZIEM DAVID DE JESUIS \and NKONGMENOCK LOIC PARFAIT \and MARTIAL JEANNOT MAA
}

\institute[UY1]{%
  Université de Yaoundé I \\
  \vspace{0.2cm}
  Faculté des Sciences, Département d'informatique
}
\date[\textcolor{white}{\today}]{Année Académique 2025-2026}

\begin{document}

\frame{\titlepage}

\begin{frame}
\frametitle{Plan}
\begin{multicols}{2}
    \begin{onehalfspace}
        \tableofcontents
    \end{onehalfspace}
\end{multicols}
\end{frame}

%------------------------------------------------------------
\section{\Large{Introduction}}
\subsection{\small Contexte, Problématique, Solution, avantages}

\begin{frame}{Contexte, Problématique, Solution et Avantages}
\begin{block}{\textbf{Contexte général}}
Au Cameroun, les services de paiement mobile (Orange Money, MTN Mobile Money) sont omniprésents mais manquent d'interopérabilité et de solutions adaptées aux commerçants. Les clients rencontrent des difficultés pour payer à la caisse rapidement et sans espèces.
\end{block}

\begin{block}{\textbf{Problématique}}
\begin{itemize}
    \item Processus de paiement lent et inefficace à la caisse
    \item Dépendance aux terminaux physiques de paiement
    \item Manque de solutions digitales intégrées pour les commerçants
    \item Frais de transaction élevés pour les petits commerces
\end{itemize}
\end{block}

\begin{block}{\textbf{Solution proposée : PayQr}}
Développement d'une plateforme complète permettant :
\begin{itemize}
    \item Génération de QR codes par les commerçants
    \item Paiement instantané par scan de QR code
    \item Gestion de portefeuille virtuel client
    \item Intégration avec les passerelles Mobile Money via Aangaraa-Pay
\end{itemize}
\end{block}

\begin{block}{\textbf{Avantages clés}}
\begin{itemize}
    \item Paiement à la caisse ultra-rapide via QR code
    \item Réduction des files d'attente en magasin
    \item Solution accessible et économique pour les commerçants
    \item Traçabilité complète des transactions
\end{itemize}
\end{block}
\end{frame}

\section{\Large Structure de stage}
\subsection{\small ARITeD}

\begin{frame}{Cadre professionnel et missions}
\begin{block}{Cadre professionnel}
    \begin{itemize}
      \item \textbf{Entreprise} : ARITeD (Association de Recherche et d'Innovation en Technologie pour le Développement)
      \item \textbf{Localisation} : Chapelle Obili, Yaoundé (près de l'UY1)
      \item \textbf{Période} : Janvier 2026 - Mai 2026
    \end{itemize}
\end{block}

\begin{block}{Missions d'ARITeD}
    \begin{itemize}
      \item Conception de logiciels sur mesure (web, mobiles) pour les entreprises
      \item Intégration de technologies favorisant la transformation numérique
      \item Incubation et accompagnement des jeunes étudiants porteurs de projets
    \end{itemize}
\end{block}
\end{frame}

\begin{frame}{Déroulement du projet}
  \begin{block}{\Large Résumé du déroulement du projet}
Le projet PayQr s'est déroulé de janvier à mai 2026 en partenariat avec ARITeD. Après une phase d'analyse des besoins des commerçants et des systèmes de paiement existants, nous avons conçu une architecture full-stack adaptée. Le développement, lancé en mars, a intégré les fonctionnalités essentielles (authentification JWT, génération QR, paiement, intégration Aangaraa) selon une méthode itérative. En mai, le projet a été finalisé avec les phases de test et de déploiement.
\end{block}
\end{frame}

\section{\Large Architecture technique}
\subsection{\large Architecture globale}

\begin{frame}{Architecture de la solution PayQr}
\begin{block}{\normalsize Architecture full-stack}
\textbf{PayQr} repose sur une architecture moderne et modulaire :
\begin{itemize}
    \item \textbf{Backend} : Spring Boot (Java) avec JWT pour l'authentification
    \item \textbf{Base de données} : MySQL 8.0 pour les transactions et utilisateurs
    \item \textbf{API externe} : Intégration Aangaraa-Pay pour Mobile Money
    \item \textbf{Frontend} : React JS pour le tableau de bord vendeur
    \item \textbf{Mobile} : Flutter pour l'application client
\end{itemize}
\end{block}

\begin{block}{\normalsize Entités principales}
\begin{itemize}
    \item \textbf{User} : Entité abstraite avec email, téléphone, mot de passe
    \item \textbf{Client} : Solde virtuel, historique transactions
    \item \textbf{Vendeur} : Nom commerce, adresse, solde virtuel caisse
    \item \textbf{QRCode} : Code unique, montant, lien au vendeur
    \item \textbf{Transaction} : Montant, statut, commission, type opération
\end{itemize}
\end{block}
\end{frame}

\subsection{\large API Backend}
\begin{frame}{Endpoints principaux - Authentification}
\begin{block}{Authentification JWT}
\begin{itemize}
    \item \texttt{POST /api/auth/register/client} - Inscription client
    \item \texttt{POST /api/auth/register/vendeur} - Inscription vendeur  
    \item \texttt{POST /api/auth/login} - Connexion avec token JWT
    \item \texttt{POST /api/auth/forgot-password} - Demande reset mot de passe
    \item \texttt{POST /api/auth/reset-password} - Réinitialisation avec code
\end{itemize}
\end{block}

\begin{block}{Gestion des comptes}
\begin{itemize}
    \item \texttt{GET /api/client/me} - Profil client authentifié
    \item \texttt{GET /api/client/solde} - Solde virtuel du client
    \item \texttt{GET /api/vendeur/me} - Profil vendeur authentifié
    \item \texttt{GET /api/vendeur/solde} - Solde caisse du vendeur
\end{itemize}
\end{block}
\end{frame}

\begin{frame}{Endpoints principaux - Paiement QR}
\begin{block}{Operations QR Code}
\begin{itemize}
    \item \texttt{POST /api/qrcode/generate} - Génération QR par le vendeur
    \item \texttt{POST /api/payments/initiate} - Initialisation paiement externe
    \item \texttt{POST /api/payments/virtual} - Paiement via solde interne
    \item \texttt{GET /api/payments/status/\{id\}} - Vérification statut transaction
    \item \texttt{POST /api/webhook/aangaraa} - Réception notifications webhook
\end{itemize}
\end{block}

\begin{block}{Retraits et solde}
\begin{itemize}
    \item \texttt{POST /api/retrait} - Demande de retrait vendeur
    \item \texttt{GET /api/transactions} - Historique des transactions
    \item \texttt{GET /api/audit-log} - Journal d'audit complet
\end{itemize}
\end{block}
\end{frame}

\section{\Large Fonctionnalités détaillées}
\subsection{\large Parcours utilisateur}

\begin{frame}{1. Création de compte et authentification}
\begin{block}{Inscription}
\begin{itemize}
    \item Client : nom, email, téléphone, mot de passe
    \item Vendeur : nom, email, téléphone, mot de passe, nom commerce, adresse
    \item Validation : email/téléphone unique, mot de passe hashé
\end{itemize}
\end{block}

\begin{block}{Authentification}
\begin{itemize}
    \item Login avec email ou téléphone + mot de passe
    \item Génération de token JWT valide 24h
    \item Rôle déterminé (CLIENT, VENDEUR, ADMIN)
    \item Endpoint \texttt{/api/auth/me} pour vérifier session active
\end{itemize}
\end{block}

\begin{block}{Mot de passe oublié}
\begin{itemize}
    \item Envoi d'un code OTP par email
    \item Validation du code et réinitialisation sécurisée
\end{itemize}
\end{block}
\end{frame}

\begin{frame}{2. Rechargement du compte virtuel}
\begin{block}{Processus de recharge}
\begin{itemize}
    \item Client saisit montant (minimum 10 XAF) et opérateur (Orange/MTN)
    \item Appel API Aangaraa-Pay pour paiement Mobile Money
    \item Retour avec payToken pour validation
    \item Webhook automatique pour crédit du solde une fois paiement validé
\end{itemize}
\end{block}

\begin{block}{Endpoints concernés}
\begin{itemize}
    \item \texttt{POST /api/payments/rechargement} (via PaymentService)
    \item \texttt{POST /api/webhook/aangaraa} - Traitement webhook
    \item \texttt{GET /api/client/solde} - Consultation solde
\end{itemize}
\end{block}

\begin{block}{Sécurité}
\begin{itemize}
    \item Transaction créée avec statut PENDING
    \item Solde crédité uniquement après confirmation webhook
    \item Scheduler de secours vérifie les transactions PENDING
\end{itemize}
\end{block}
\end{frame}

\begin{frame}{3. Génération et paiement QR Code}
\begin{block}{Par le vendeur}
\begin{itemize}
    \item Génération d'un QR code avec montant et description
    \item QR code lié au vendeur (compte caisse)
    \item Expiration configurable (30 minutes par défaut)
\end{itemize}
\end{block}

\begin{block}{Par le client}
\begin{itemize}
    \item Scan du QR code avec l'application mobile
    \item Option 1 : Paiement externe (Mobile Money via Aangaraa)
    \item Option 2 : Paiement interne (solde virtuel)
    \item Validation instantanée et confirmation
\end{itemize}
\end{block}

\begin{block}{Traitement backend}
\begin{itemize}
    \item Vérification expiration et statut utilisation
    \item Calcul commission (configurable)
    \item Mise à jour solde vendeur automatique
    \item Statut QR changé à "utilisé" après paiement
\end{itemize}
\end{block}
\end{frame}

\begin{frame}{4. Retrait des fonds}
\begin{block}{Processus de retrait}
\begin{itemize}
    \item Vendeur demande retrait vers son Mobile Money
    \item Choix opérateur (Orange Cameroon / MTN Cameroon)
    \item Appel API Aangaraa pour traitement du retrait
    \item Webhook de confirmation et débit automatique du solde
\end{itemize}
\end{block}

\begin{block}{Endpoints concernés}
\begin{itemize}
    \item \texttt{POST /api/retrait} - Création demande retrait
    \item \texttt{GET /api/retrait/history} - Historique retraits
    \item Scheduler \texttt{checkPendingRetraits()} - Vérification automatique
\end{itemize}
\end{block}

\begin{block}{Gestion des erreurs}
\begin{itemize}
    \item Retraits expirés (> 15 min) marqués FAILED
    \item Retry automatique via scheduler 30 secondes
    \item Notification des erreurs en log
\end{itemize}
\end{block}
\end{frame}

\section{\Large Implémentation détaillée}
\subsection{\large Services clés}

\begin{frame}{PaymentService - Cœur métier}
\begin{block}{Fonctions principales}
\begin{itemize}
    \item \textbf{initierRechargement()} : Crée transaction RECHARGEMENT et appelle Aangaraa
    \item \textbf{initiatePayment()} : Gère paiements QR externes avec commission
    \item \textbf{initiateVirtualPayment()} : Paiement interne solde client $\rightarrow$ vendeur
    \item \textbf{checkPaymentStatus()} : Vérification statut auprès d'Aangaraa
    \item \textbf{processWebhook()} : Traitement notifications automatiques
\end{itemize}
\end{block}

\begin{block}{Schedulers automatiques}
\begin{itemize}
    \item \texttt{checkPendingTransactions()} : Toutes les 30s, vérifie PENDING
    \item \texttt{checkPendingRetraits()} : Toutes les 30s, suit les retraits
    \item Expiration automatique après 15 minutes d'inactivité
    \item Fallback si webhook non reçu
\end{itemize}
\end{block}
\end{frame}

\begin{frame}{Sécurité et authentification}
\begin{block}{JwtUtil - Génération des tokens}
\begin{itemize}
    \item Algorithm HS256 avec secret partagé
    \item Claims : username, userId, role, expiration
    \item Token signé et renvoyé au client après login
\end{itemize}
\end{block}

\begin{block}{JwtAuthenticationFilter}
\begin{itemize}
    \item Intercepte chaque requête HTTP
    \item Extrait et valide le token JWT
    \item Charge l'utilisateur dans le contexte Spring Security
    \item Endpoints publics : /api/auth/*, /api/webhook/*
    \item Endpoints protégés : nécessitent authentification
\end{itemize}
\end{block}

\begin{block}{PasswordEncoderUtil}
\begin{itemize}
    \item Hashage BCrypt des mots de passe
    \item Validation lors de la connexion
\end{itemize}
\end{block}
\end{frame}

\begin{frame}{Gestion des transactions}
\begin{block}{Entité Transaction}
\begin{itemize}
    \item Statuts : PENDING, SUCCESSFUL, FAILED, CANCELLED
    \item Types : RECHARGEMENT, PAYMENT\_MARCHAND, TRANSFERT\_VIRTUEL
    \item Champs : montant, commission, montantNet, payToken, referenceOperateur
    \item Liens : client, vendeur (via QR), QRCode
\end{itemize}
\end{block}

\begin{block}{Calcul des commissions}
\begin{itemize}
    \item Configuration via \texttt{ConfigurationFrais} (taux configurable)
    \item Commission déduite du montant brut
    \item Le vendeur reçoit le montant net
    \item Exemple : 1000 XAF avec 2\% $\rightarrow$ 980 XAF au vendeur
\end{itemize}
\end{block}
\end{frame}

\section{\Large Tests \& Résultats}
\subsection{\large Résultats techniques}

\begin{frame}{Résultats techniques observés}
\begin{block}{Temps moyen par fonctionnalité}
\begin{footnotesize}
\begin{tabular}{|p{6cm}|p{3cm}|}
\hline
\rowcolor{myNewColorA!40} \textbf{Fonctionnalité} & \textbf{Temps moyen (ms)} \\
\hline
Authentification et génération de token JWT & 540 \\
Rechargement compte virtuel (Mobile Money) & 850 \\
Paiement QR externe (Aangaraa + webhook) & 1200 \\
Paiement QR interne (solde virtuel) & 670 \\
Génération QR Code & 200 \\
Retrait vers Mobile Money & 900 \\
Consultation du solde & 250 \\
Chargement historique transactions & 400 \\
\hline
\end{tabular}
\end{footnotesize}
\end{block}

\begin{block}{Taux de réussite}
94 opérations réussies sur 100 simulées : \textbf{94\%} de succès \\
\textit{Causes d'échec : coupures réseau, erreurs côté Aangaraa.}
\end{block}
\end{frame}

\begin{frame}{Points forts et limites}
\begin{block}{Points forts observés}
\begin{itemize}
    \item Architecture modulaire et extensible
    \item Intégration Aangaraa-Pay robuste avec fallbacks
    \item Paiement instantané avec solde virtuel
    \item Sécurité JWT renforcée
    \item Webhooks automatiques avec scheduler de secours
\end{itemize}
\end{block}

\begin{block}{Limites identifiées}
\begin{itemize}
    \item Dépendance à l'API Aangaraa pour Mobile Money
    \item Pas de support iOS natif (Flutter en cours)
    \item Gestion des erreurs réseau à améliorer
    \item Nécessite validation réglementaire pour production
\end{itemize}
\end{block}
\end{frame}

\section{\Large Conclusion}
\begin{frame}
\begin{onehalfspace}
\begin{block}{\Large Bilan général}
Les solutions de paiement mobile décentralisées offrent une alternative innovante aux systèmes traditionnels. \textbf{PayQr} répond à un besoin local concret : offrir aux commerçants une interface de caisse unifiée, fluide et sécurisée.

\vspace{0.2cm}
Notre solution, reposant sur Spring Boot, MySQL et l'intégration Aangaraa-Pay, permet recharges, paiements par QR code et retraits Mobile Money de façon fluide et simple.
\end{block}

\begin{block}{\Large Perspectives}
\begin{itemize}
    \item Déploiement en phase pilote avec commerçants locaux
    \item Intégration de nouvelles passerelles de paiement
    \item Développement du portail web vendeur React complet
    \item Obtention des autorisations nécessaires (MINMAP, ART)
    \item Publication sur Play Store et App Store
\end{itemize}
\end{block}
\end{onehalfspace}
\end{frame}

\section{\large Démonstration}
\begin{frame}
\begin{center}
\includegraphics[width=0.8\textwidth]{ressources/images/demo.png}
\captionof{figure}{Interface de paiement QR Code - PayQr}
\end{center}
\end{frame}

\section*{\large Remerciements}   
\begin{frame}
\textcolor{myNewColorA}{\Huge{\centerline{Merci pour votre aimable attention!}}}

\vspace{1cm}
\begin{center}
\textbf{Questions \& Réponses}
\end{center}
\end{frame}

\end{document}